\documentclass[10pt,a4paper]{amsart}
\usepackage[margin=19mm]{geometry}
\usepackage{amsmath,amssymb,mathtools}
\usepackage[scr=rsfs,cal=euler]{mathalfa}
\usepackage{xfrac}
\usepackage[unicode,hidelinks,bookmarks=false]{hyperref}
\newcommand{\ie}{i.e.\ }
\newcommand{\cf}{cf.\ }
\newcommand{\resp}{respectively\ }
\newcommand{\Subsection}{\subsection}
\providecommand{\nobreakdash}{\nobreak\hbox{-}}
\setlength{\emergencystretch}{2em}
\multlinegap0pt
\usepackage{amssymb}
\usepackage[shortlabels]{enumitem}
\usepackage{float}
\usepackage{mathtools}
\usepackage{cancel}
\newtheorem{theorem}{Theorem}
\newtheorem{lemma}{Lemma}
\newcommand{\Peri}{\operatorname{Peri}}
\newcommand{\bb}{\mathbb}
\title{Topological Data Analysis Meets Convex Geometry: New Perspectives on the Euler Characteristic Transform}
\author{Jes\'us Gac\'ias Franco}
\date{Source: arXiv:2607.28021v1 (CC BY 4.0)}
\begin{document}
\maketitle

\noindent\emph{Source and licence.} Topological Data Analysis Meets Convex Geometry: New Perspectives on the Euler Characteristic Transform. Version arXiv:2607.28021v1 (CC BY 4.0), distributed by arXiv under the Creative Commons Attribution 4.0 International licence, http://creativecommons.org/licenses/by/4.0/, as recorded in the arXiv licence metadata for this exact version and shown on the abstract page https://arxiv.org/abs/2607.28021v1. Full licence notice: \texttt{licenses/arxiv-2607.28021-CC-BY-4.0.txt}.\par\smallskip

\noindent\textit{Editorial scope.} Counted unit: the article's theorem Approx (source0.tex lines 1333--1342). The article's complete original proof of it is delivered with it (source0.tex lines 1360--1389, one \texttt{proof} environment of 30 lines, which the article places away from the statement). No mathematics is written, repaired or re-derived: the statement and the proof are the article's own lines, in the article's own notation, and no retained line is substituted or re-flowed.  Retained uncounted as the source-local context the proof needs: lines 1344--1347.  One declared adaptation: exactly 1 ancillary strings inside the retained spans are omitted (image-inclusion commands and, where the source repeats a label, the repeated label definitions) and no other line differs from the source; the figure environments, with their placement, captions and labels, are kept, and the package records the omitted strings in fidelity receipts bound to the affected bodies.  No retained line contains a citation, so the wrapper carries no bibliography and no bibliography entry.  The retained lines contain the article's own diagram code, which the wrapper typesets with the standard tikz packages; no image file is delivered and the package declares no asset dependency.  Editorial formatting only, in addition: an AMS \texttt{amsart} preamble that restates the article's own declarations and macros used by the retained lines, namely the environments \texttt{theorem}, \texttt{lemma} on separate counters, together with the standard packages those lines need.  The retained environments print in this document as Theorem 1, Lemma 1 in order of appearance, whereas the article prints the same items with the numbering of its own full document.  The complete native source of the article as distributed in the arXiv e-print for this exact version is bundled unchanged as \texttt{sources/206382/source0.tex}.
\begin{lemma}
	\label{ConvexCurve}
	Let $C \subseteq \bb R^2$ be a convex, compact, tame, and $\mathcal C^1$ curve joining two horizontally aligned points. Then, a sequence $\{S_n\}_{n \in \bb N}$ as in theorem \ref{Approx} exists. In particular, the perimeter recovery formula holds for $C$.
\end{lemma}

\begin{theorem}
	\label{Approx}
	Let $M \subseteq \bb R^2$ be a compact and tame shape. There exists a sequence of simplicial complexes $\{S_n\}_{n \in \bb N}$ such that:
	
	\begin{enumerate}
		\item $S_n$ is homeomorphic to $M$ for all $n$.
		\item $\Peri(S_n) \overset{n \to \infty}{\longrightarrow} \Peri(M)$.
		\item $\int_{\bb S^{1}}h_{S_n}(v) dv \overset{n \to \infty}{\longrightarrow} \int_{\bb S^{1}} h_M(v) dv$.
	\end{enumerate}
\end{theorem}

\begin{proof}[Proof of theorem \ref{Approx}]
	We first describe a triangulation of $\partial M$ (that is, the image of boundary and isolated edges of any triangulation of $M$). Since $\partial M$ consists of finitely many $\mathcal C^1$ curves, consider for each an approximation by polygonal curves whose vertices lie in $\partial M$. Choose these decompositions fine enough so that for each line segment of the polygonal curves, the arc of $\partial M$ drawn between the two points it joins is either concave, convex, or perfectly matches the line segment. Naming $S$ this set of polygonal curves, let $f: S \to \partial M$ be this $\mathcal C^1$ triangulation. Note that, by construction, all vertices of $S$ are fixed points of $f$. We now extend $f$ to a triangulation $\tilde f$ of $M$, as follows:
	\begin{itemize}
		\item Triangulate the convex closure of $S$ in a compatible way with the natural triangulation of $S$. By applying barycentric subdivision if necessary, suppose each added face contains at most one edge of $S$ on its boundary.
		\item Fix $\tilde f(v) = v$ for all vertices $v$ of the resulting triangulation, and $\tilde f|_{S}=f$. Extend $\tilde f$ linearly to all added edges, and then to all faces (non-linearly on faces with a boundary edge).
		\item The interior of each $2$-simplex now either maps completely inside $M$ or outside $M$. Only add to $S$ all those who map to $M$, together with their boundaries.
		\item Finally, add any isolated points $M$ might have to $S$ and extend $\tilde f$ as the identity on them.
	\end{itemize}
	
	By repeating this process with finer and finer triangulations of $\partial M$, we obtain a sequence $S_n$ of simplicial complexes who, by construction, verify the first two points of the theorem. As for the third, let us show that
	\[ \int_{\bb S^{1}} (h_{S_n}(v)-h_{M}(v))  dv \overset{n \to \infty} \longrightarrow 0\, .\]
	Using additivity, decompose $S_n$ as the union of its faces, isolated edges, and isolated points, and decompose $M$ as the image under $\tilde f$ of that same union. By construction of $\tilde f$, faces that don't contain boundary edges, internal edges and all vertices remain fixed under $\tilde f$. These terms cancel out, so the integral equals
	\[ \sum_{F} \int(h_F -h_{\tilde f(F)}) + \sum_{e} \int(h_e -h_{\tilde f(e)})\, , \]
	where the first sum ranges over faces of $S_n$ with a boundary edge, and the second over isolated edges. Focusing first on the isolated edges, $\int h_e$ is twice the length of the straight-line segment $e$, and $h_{\tilde f(e)}$ is twice the length of the arc of $M$ between the endpoints of $e$, by lemma \ref{ConvexCurve}.
	
	For faces, we distinguish between two cases: for a face $F$, $\tilde f(F)$ differs from $F$ by either introducing a ``bump'' or a ``dent'' (see figure \ref{BumpDent}), depending on whether the image under $\tilde f $ of its boundary edge is concave or convex.
	\begin{figure}
		\centering
	
	\caption{Two faces of the triangulation $S_n$ and their images under $\tilde f$. The first is a bump, the second is a dent. In each case, the part added or removed to pass from $F$ to $\tilde f(F)$ is convex.}
	\label{BumpDent}
	\end{figure}
	In the case of a bump, we have $\tilde f(F)= F \cup B$, with $B$ being the added, convex, bump. Thus, by additivity $h_F-h_{\tilde f(F)}= h_{F \cap B}-h_B$. Now, $F \cap B$ is nothing but the boundary edge of $F$, so both terms on the right are convex. Since the perimeter formula holds for convex shapes\footnote{If one wishes to not use this fact, a similar proof to the one of lemma \ref{ConvexCurve} shows the assertion for bumps.}, the difference of their integrals turns out to be the length of the boundary edge minus the length of its image under $\tilde f$.
	
	In the case of dent, write $F= \tilde f(F) \cup D$, with $D= \overline{F \setminus \tilde f(F)}$ being the dent, and applying additivity, we have $f_F-h_{\tilde f(F)}= h_D-h_{D \cap \tilde f(F)}$. $D$ is again convex, with one component of its boundary being the boundary edge of $F$, and the other its image under $\tilde f$. On the other hand, $D \cap \tilde f(F)$ is a convex curve, so by applying lemma \ref{ConvexCurve}, the difference of integrals again turns out to be the length of the boundary edge minus the length of its image under $\tilde f$.
	
	Collecting terms, the difference of integrals for any fixed $n$ equals:
	\[\Peri(S_n)-\Peri(M). \]
	Since $S_n$ verifies point two, this converges to zero.
	\end{proof}

\end{document}
